\subsubsection{Factor hunting}

The next trick -- known as \emph{factor hunting} -- works not only for
determinants; however, determinants provide some of the simplest examples.

\begin{theorem}
[Vandermonde determinant]\label{thm.det.vander}Let $n\in\mathbb{N}$. Let
$a_{1},a_{2},\ldots,a_{n}$ be $n$ elements of $K$. Then: \medskip

\textbf{(a)} We have%
\[
\det\left(  \left(  a_{i}^{n-j}\right)  _{1\leq i\leq n,\ 1\leq j\leq
n}\right)  =\prod_{1\leq i<j\leq n}\left(  a_{i}-a_{j}\right)  .
\]


\textbf{(b)} We have%
\[
\det\left(  \left(  a_{j}^{n-i}\right)  _{1\leq i\leq n,\ 1\leq j\leq
n}\right)  =\prod_{1\leq i<j\leq n}\left(  a_{i}-a_{j}\right)  .
\]


\textbf{(c)} We have%
\[
\det\left(  \left(  a_{i}^{j-1}\right)  _{1\leq i\leq n,\ 1\leq j\leq
n}\right)  =\prod_{1\leq j<i\leq n}\left(  a_{i}-a_{j}\right)  .
\]


\textbf{(d)} We have%
\[
\det\left(  \left(  a_{j}^{i-1}\right)  _{1\leq i\leq n,\ 1\leq j\leq
n}\right)  =\prod_{1\leq j<i\leq n}\left(  a_{i}-a_{j}\right)  .
\]

\end{theorem}

\begin{example}
Here is what the four parts of Theorem \ref{thm.det.vander} say for $n=3$:
\medskip

\textbf{(a)} We have $\det\left(
\begin{array}
[c]{ccc}%
a_{1}^{2} & a_{1} & 1\\
a_{2}^{2} & a_{2} & 1\\
a_{3}^{2} & a_{3} & 1
\end{array}
\right)  =\left(  a_{1}-a_{2}\right)  \left(  a_{1}-a_{3}\right)  \left(
a_{2}-a_{3}\right)  $. \medskip

\textbf{(b)} We have $\det\left(
\begin{array}
[c]{ccc}%
a_{1}^{2} & a_{2}^{2} & a_{3}^{2}\\
a_{1} & a_{2} & a_{3}\\
1 & 1 & 1
\end{array}
\right)  =\left(  a_{1}-a_{2}\right)  \left(  a_{1}-a_{3}\right)  \left(
a_{2}-a_{3}\right)  $. \medskip

\textbf{(c)} We have $\det\left(
\begin{array}
[c]{ccc}%
1 & a_{1} & a_{1}^{2}\\
1 & a_{2} & a_{2}^{2}\\
1 & a_{3} & a_{3}^{2}%
\end{array}
\right)  =\left(  a_{2}-a_{1}\right)  \left(  a_{3}-a_{1}\right)  \left(
a_{3}-a_{2}\right)  $. \medskip

\textbf{(d)} We have $\det\left(
\begin{array}
[c]{ccc}%
1 & 1 & 1\\
a_{1} & a_{2} & a_{3}\\
a_{1}^{2} & a_{2}^{2} & a_{3}^{2}%
\end{array}
\right)  =\left(  a_{2}-a_{1}\right)  \left(  a_{3}-a_{1}\right)  \left(
a_{3}-a_{2}\right)  $.
\end{example}

Theorem \ref{thm.det.vander} is a classical and important result, known as the
\emph{Vandermonde determinant}. Many different proofs are known (see, e.g.,
\cite[Theorem 6.46]{detnotes} or \cite[\S 5.3]{Aigner07} or \cite[Section
III.8.6, Example (1)]{Bourba74} or \cite[Theorem 1]{GriHyp}; a combinatorial
proof can also be found in Exercise \ref{exe.det.vander-tour}; two more proofs
are obtained in Exercise \ref{exe.det.vander.kratt-lem6} and Exercise
\ref{exe.det.binom-by-vander}). We will now sketch a proof using factor
hunting and polynomials. We will first focus on proving part \textbf{(a)} of
Theorem \ref{thm.det.vander}, and afterwards derive the other parts from it.

The first step of our proof is reducing Theorem \ref{thm.det.vander}
\textbf{(a)} to the \textquotedblleft particular case\textquotedblright\ in
which $K$ is the polynomial ring $\mathbb{Z}\left[  x_{1},x_{2},\ldots
,x_{n}\right]  $ and the elements $a_{1},a_{2},\ldots,a_{n}$ are the
indeterminates $x_{1},x_{2},\ldots,x_{n}$. This is merely a particular case
(one possible choice of $K$ and $a_{1},a_{2},\ldots,a_{n}$ among many);
however, as we will soon see, proving Theorem \ref{thm.det.vander}
\textbf{(a)} in this particular case will quickly entail that Theorem
\ref{thm.det.vander} \textbf{(a)} holds in the general case. Let us elaborate
on this argument. First, let us state Theorem \ref{thm.det.vander}
\textbf{(a)} in this particular case as a lemma:

\begin{lemma}
\label{lem.det.vander.a.pol}Let $n\in\mathbb{N}$. Consider the polynomial ring
$\mathbb{Z}\left[  x_{1},x_{2},\ldots,x_{n}\right]  $ in $n$ indeterminates
$x_{1},x_{2},\ldots,x_{n}$ with integer coefficients. In this ring, we have%
\begin{equation}
\det\left(  \left(  x_{i}^{n-j}\right)  _{1\leq i\leq n,\ 1\leq j\leq
n}\right)  =\prod_{1\leq i<j\leq n}\left(  x_{i}-x_{j}\right)  .
\label{eq.lem.det.vander.a.pol.1}%
\end{equation}

\end{lemma}

We can derive Theorem \ref{thm.det.vander} \textbf{(a)} from Lemma
\ref{lem.det.vander.a.pol} as follows:

\begin{proof}
[Proof of Theorem \ref{thm.det.vander} \textbf{(a)} using Lemma
\ref{lem.det.vander.a.pol}.]The equality%
\begin{equation}
\det\left(  \left(  a_{i}^{n-j}\right)  _{1\leq i\leq n,\ 1\leq j\leq
n}\right)  =\prod_{1\leq i<j\leq n}\left(  a_{i}-a_{j}\right)
\label{pf.thm.det.vander.a.1}%
\end{equation}
follows from the equality (\ref{eq.lem.det.vander.a.pol.1}) by substituting
$a_{1},a_{2},\ldots,a_{n}$ for $x_{1},x_{2},\ldots,x_{n}$. \medskip

\begin{fineprint}
This is sufficiently clear to be considered a complete proof, but just in
case, here are a few details.

We can substitute $a_{1},a_{2},\ldots,a_{n}$ for $x_{1},x_{2},\ldots,x_{n}$ in
any polynomial $f\in\mathbb{Z}\left[  x_{1},x_{2},\ldots,x_{n}\right]  $,
since $a_{1},a_{2},\ldots,a_{n}$ are $n$ elements of a commutative ring
(namely, of $K$). It is obvious that $\prod_{1\leq i<j\leq n}\left(
x_{i}-x_{j}\right)  $ becomes $\prod_{1\leq i<j\leq n}\left(  a_{i}%
-a_{j}\right)  $ when we substitute $a_{1},a_{2},\ldots,a_{n}$ for
$x_{1},x_{2},\ldots,x_{n}$. It is perhaps a bit less obvious that $\det\left(
\left(  x_{i}^{n-j}\right)  _{1\leq i\leq n,\ 1\leq j\leq n}\right)  $ becomes
$\det\left(  \left(  a_{i}^{n-j}\right)  _{1\leq i\leq n,\ 1\leq j\leq
n}\right)  $ when we substitute $a_{1},a_{2},\ldots,a_{n}$ for $x_{1}%
,x_{2},\ldots,x_{n}$. To convince our skeptical selves of this, we expand both
determinants: The definition of the determinant yields%
\begin{align*}
\det\left(  \left(  x_{i}^{n-j}\right)  _{1\leq i\leq n,\ 1\leq j\leq
n}\right)   &  =\sum_{\sigma\in S_{n}}\left(  -1\right)  ^{\sigma}%
x_{1}^{n-\sigma\left(  1\right)  }x_{2}^{n-\sigma\left(  2\right)  }\cdots
x_{n}^{n-\sigma\left(  n\right)  }\ \ \ \ \ \ \ \ \ \ \text{and}\\
\det\left(  \left(  a_{i}^{n-j}\right)  _{1\leq i\leq n,\ 1\leq j\leq
n}\right)   &  =\sum_{\sigma\in S_{n}}\left(  -1\right)  ^{\sigma}%
a_{1}^{n-\sigma\left(  1\right)  }a_{2}^{n-\sigma\left(  2\right)  }\cdots
a_{n}^{n-\sigma\left(  n\right)  },
\end{align*}
and it is clear that substituting $a_{1},a_{2},\ldots,a_{n}$ for $x_{1}%
,x_{2},\ldots,x_{n}$ transforms \newline$\sum_{\sigma\in S_{n}}\left(
-1\right)  ^{\sigma}x_{1}^{n-\sigma\left(  1\right)  }x_{2}^{n-\sigma\left(
2\right)  }\cdots x_{n}^{n-\sigma\left(  n\right)  }$ into $\sum_{\sigma\in
S_{n}}\left(  -1\right)  ^{\sigma}a_{1}^{n-\sigma\left(  1\right)  }%
a_{2}^{n-\sigma\left(  2\right)  }\cdots a_{n}^{n-\sigma\left(  n\right)  }$.
\end{fineprint}

Thus, (\ref{pf.thm.det.vander.a.1}) follows from
(\ref{eq.lem.det.vander.a.pol.1}). In other words, Theorem
\ref{thm.det.vander} \textbf{(a)} follows from Lemma
\ref{lem.det.vander.a.pol}.
\end{proof}

Arguments like the one we just used are frequently applied in algebra; see
\cite{Conrad-ui} for some more examples.

It now remains to prove Lemma \ref{lem.det.vander.a.pol}.

\begin{proof}
[Proof of Lemma \ref{lem.det.vander.a.pol} (sketched).]We set%
\[
f:=\det\left(  \left(  x_{i}^{n-j}\right)  _{1\leq i\leq n,\ 1\leq j\leq
n}\right)  \ \ \ \ \ \ \ \ \ \ \text{and}\ \ \ \ \ \ \ \ \ \ g:=\prod_{1\leq
i<j\leq n}\left(  x_{i}-x_{j}\right)  .
\]
Thus, we must prove that $f=g$.

We have%
\begin{align}
f  &  =\det\left(  \left(  x_{i}^{n-j}\right)  _{1\leq i\leq n,\ 1\leq j\leq
n}\right) \nonumber\\
&  =\sum_{\sigma\in S_{n}}\left(  -1\right)  ^{\sigma}x_{1}^{n-\sigma\left(
1\right)  }x_{2}^{n-\sigma\left(  2\right)  }\cdots x_{n}^{n-\sigma\left(
n\right)  } \label{pf.lem.det.vander.a.pol.f=1}%
\end{align}
(by the definition of a determinant). The right hand side of this equality is
a homogeneous polynomial in $x_{1},x_{2},\ldots,x_{n}$ of degree
$\dfrac{n\left(  n-1\right)  }{2}$\ \ \ \ \footnote{because it is a
$\mathbb{Z}$-linear combination of the monomials $x_{1}^{n-\sigma\left(
1\right)  }x_{2}^{n-\sigma\left(  2\right)  }\cdots x_{n}^{n-\sigma\left(
n\right)  }$, each of which has degree%
\begin{align*}
\left(  n-\sigma\left(  1\right)  \right)  +\left(  n-\sigma\left(  2\right)
\right)  +\cdots+\left(  n-\sigma\left(  n\right)  \right)   &  =n\cdot
n-\underbrace{\left(  \sigma\left(  1\right)  +\sigma\left(  2\right)
+\cdots+\sigma\left(  n\right)  \right)  }_{\substack{=1+2+\cdots
+n\\\text{(since }\sigma\text{ is a permutation of }\left[  n\right]
\text{)}}}\\
&  =n\cdot n-\underbrace{\left(  1+2+\cdots+n\right)  }_{=\dfrac{n\left(
n+1\right)  }{2}}=n\cdot n-\dfrac{n\left(  n+1\right)  }{2}\\
&  =\dfrac{n\left(  n-1\right)  }{2}%
\end{align*}
}. Thus, $f$ is a homogeneous polynomial in $x_{1},x_{2},\ldots,x_{n}$ of
degree $\dfrac{n\left(  n-1\right)  }{2}$. Furthermore, the monomial
$x_{1}^{n-1}x_{2}^{n-2}\cdots x_{n}^{n-n}$ appears with coefficient $1$ on the
right hand side of (\ref{pf.lem.det.vander.a.pol.f=1}) (indeed, all the $n!$
addends in the sum on this right hand side contain distinct monomials, and
thus only the addend for $\sigma=\operatorname*{id}$ makes any contribution to
the coefficient of the monomial $x_{1}^{n-1}x_{2}^{n-2}\cdots x_{n}^{n-n}$).
Hence,%
\begin{equation}
\left[  x_{1}^{n-1}x_{2}^{n-2}\cdots x_{n}^{n-n}\right]  f=1.
\label{pf.lem.det.vander.a.pol.f-LC}%
\end{equation}
Therefore, $f\neq0$.

Now, let $u$ and $v$ be two elements of $\left[  n\right]  $ satisfying $u<v$.
Then, the polynomial $f$ becomes $0$ when we set $x_{u}$ equal to $x_{v}$
(that is, when we substitute $x_{v}$ for $x_{u}$). Indeed, when we set $x_{u}$
equal to $x_{v}$, the matrix $\left(  x_{i}^{n-j}\right)  _{1\leq i\leq
n,\ 1\leq j\leq n}$ becomes a matrix that has two equal rows (namely, its
$u$-th and its $v$-th row both become equal to $\left(  x_{v}^{n-1}%
,x_{v}^{n-2},\ldots,x_{v}^{n-n}\right)  $), and thus its determinant becomes
$0$ (by Theorem \ref{thm.det.rowop} \textbf{(c)}); but this means precisely
that $f=0$ (since $f$ is the determinant of this matrix).

Now, we recall the following well-known property of univariate polynomials:

\begin{statement}
\textit{Root factoring-off theorem:} Let $R$ be a commutative ring. Let $p\in
R\left[  t\right]  $ be a univariate polynomial that has a root $r\in R$.
Then, this polynomial $p$ is divisible by $t-r$ (in the ring $R\left[
t\right]  $).
\end{statement}

Using this property, we can easily see that if a polynomial $p\in
\mathbb{Z}\left[  x_{1},x_{2},\ldots,x_{n}\right]  $ becomes $0$ when we set
$x_{u}$ equal to $x_{v}$, then this polynomial $p$ is a multiple of
$x_{u}-x_{v}$ (in the ring $\mathbb{Z}\left[  x_{1},x_{2},\ldots,x_{n}\right]
$). (Indeed, viewing $p$ as a polynomial in the single indeterminate $x_{u}$
over the ring $\mathbb{Z}\left[  x_{1},x_{2},\ldots,\widehat{x_{u}}%
,\ldots,x_{n}\right]  $\ \ \ \ \footnote{The meaning of the hat over the
\textquotedblleft$x_{u}$\textquotedblright\ is as in Proposition
\ref{prop.det.x+ai}: It signifies that the entry $x_{u}$ is omitted from the
list (that is, \textquotedblleft$x_{1},x_{2},\ldots,\widehat{x_{u}}%
,\ldots,x_{n}$\textquotedblright\ means \textquotedblleft$x_{1},x_{2}%
,\ldots,x_{u-1},x_{u+1},x_{u+2},\ldots,x_{n}$\textquotedblright).}, we see
that $x_{v}$ is a root of $p$, and therefore the root factoring-off theorem
yields that $p$ is divisible by $x_{u}-x_{v}$.\ \ \ \ \footnote{Here, we are
tacitly using the canonical isomorphism between the polynomial rings
\[
\mathbb{Z}\left[  x_{1},x_{2},\ldots,x_{n}\right]
\ \ \ \ \ \ \ \ \ \ \text{and}\ \ \ \ \ \ \ \ \ \ \left(  \mathbb{Z}\left[
x_{1},x_{2},\ldots,\widehat{x_{u}},\ldots,x_{n}\right]  \right)  \left[
x_{u}\right]  .
\]
This isomorphism allows us to treat multivariate polynomials in $\mathbb{Z}%
\left[  x_{1},x_{2},\ldots,x_{n}\right]  $ as univariate polynomials in the
indeterminate $x_{u}$ over the ring $\mathbb{Z}\left[  x_{1},x_{2}%
,\ldots,\widehat{x_{u}},\ldots,x_{n}\right]  $, and vice versa (and ensures
that the notion of divisibility does not change between the former and the
latter).})

Applying this observation to $p=f$, we conclude that $f$ is a multiple of
$x_{u}-x_{v}$ (in the ring $\mathbb{Z}\left[  x_{1},x_{2},\ldots,x_{n}\right]
$), since we know that $f$ becomes $0$ when we set $x_{u}$ equal to $x_{v}$.

Forget that we fixed $u$ and $v$. We thus have shown that $f$ is a multiple of
$x_{u}-x_{v}$ (in the ring $\mathbb{Z}\left[  x_{1},x_{2},\ldots,x_{n}\right]
$) whenever $u$ and $v$ are two elements of $\left[  n\right]  $ satisfying
$u<v$. In other words, $f$ is a multiple of each of the linear polynomials%
\[%
\begin{array}
[c]{cccc}%
x_{1}-x_{2}, & x_{1}-x_{3}, & \ldots, & x_{1}-x_{n},\\
& x_{2}-x_{3}, & \ldots, & x_{2}-x_{n},\\
&  & \ddots & \vdots\\
&  &  & x_{n-1}-x_{n}%
\end{array}
\]
(in the ring $\mathbb{Z}\left[  x_{1},x_{2},\ldots,x_{n}\right]  $). However,
these linear polynomials are irreducible\footnote{Check this! (Actually, any
linear polynomial over $\mathbb{Z}$ is irreducible if the gcd of its
coefficients is $1$.)} and mutually non-associate (i.e., no two of them are
associate\footnote{Recall that two elements $a$ and $b$ of a commutative ring
$R$ are said to be \emph{associate} if there exists some unit $u$ of $R$ such
that $a=ub$. Being associate is known (and easily verified) to be an
equivalence relation.})\footnote{Check this! (Recall that the only units of
the polynomial ring $\mathbb{Z}\left[  x_{1},x_{2},\ldots,x_{n}\right]  $ are
$1$ and $-1$.)}. Since the ring $\mathbb{Z}\left[  x_{1},x_{2},\ldots
,x_{n}\right]  $ is a unique factorization domain\footnote{This is a
nontrivial, but rather well-known result in abstract algebra. Proofs can be
found in \cite[Theorem 3.7.4]{Ford21}, \cite[Remark after Corollary
8.21]{Knapp16}, \cite[Chapter IV, Theorems 4.8 and 4.9]{MiRiRu87} and
\cite[Essay 1.4, Corollary of Theorem 1 and Corollary 1 of Theorem
2]{Edward22}.}, we thus conclude that any polynomial $p\in\mathbb{Z}\left[
x_{1},x_{2},\ldots,x_{n}\right]  $ that is a multiple of each of these linear
polynomials must necessarily be a multiple of their product $\prod_{1\leq
i<j\leq n}\left(  x_{i}-x_{j}\right)  $. Hence, the polynomial $f$ must be a
multiple of this product $\prod_{1\leq i<j\leq n}\left(  x_{i}-x_{j}\right)  $
(since $f$ is a multiple of each of the linear polynomials above). In other
words, the polynomial $f$ must be a multiple of $g$ (since $g=\prod_{1\leq
i<j\leq n}\left(  x_{i}-x_{j}\right)  $). In other words, $f=gq$ for some
$q\in\mathbb{Z}\left[  x_{1},x_{2},\ldots,x_{n}\right]  $. Consider this $q$.
Clearly, $gq=f\neq0$ and thus $q\neq0$ and $g\neq0$.

The ring $\mathbb{Z}$ is an integral domain. Thus, any two nonzero polynomials
$a$ and $b$ over $\mathbb{Z}$ satisfy $\deg\left(  ab\right)  =\deg a+\deg b$.
Applying this to $a=g$ and $b=q$, we find $\deg\left(  gq\right)  =\deg g+\deg
q$. In other words, $\deg f=\deg g+\deg q$ (since $gq=f$).

Now, $g=\prod_{1\leq i<j\leq n}\left(  x_{i}-x_{j}\right)  $ and thus%
\begin{align*}
\deg g  &  =\deg\left(  \prod_{1\leq i<j\leq n}\left(  x_{i}-x_{j}\right)
\right)  =\sum_{1\leq i<j\leq n}\underbrace{\deg\left(  x_{i}-x_{j}\right)
}_{=1}=\sum_{1\leq i<j\leq n}1\\
&  =\left(  \text{\# of pairs }\left(  i,j\right)  \in\left[  n\right]
^{2}\text{ satisfying }i<j\right)  =\dbinom{n}{2}=\dfrac{n\left(  n-1\right)
}{2}.
\end{align*}
However, we have $\deg f\leq\dfrac{n\left(  n-1\right)  }{2}$ (since $f$ is a
homogeneous polynomial of degree $\dfrac{n\left(  n-1\right)  }{2}$). In view
of $\deg g=\dfrac{n\left(  n-1\right)  }{2}$, this rewrites as $\deg f\leq\deg
g$. Hence,%
\[
\deg g\geq\deg f=\deg g+\deg q.
\]
Therefore, $0\geq\deg q$. This shows that the polynomial $q$ is constant. In
other words, $q\in\mathbb{Z}$.

It remains to show that this constant $q$ is $1$. However, this can be done by
comparing some coefficients of $f$ and $g$. Indeed, let us look at the
coefficient of $x_{1}^{n-1}x_{2}^{n-2}\cdots x_{n}^{n-n}$ (as we already know
this coefficient for $f$ to be $1$).

Expanding the product $\prod_{1\leq i<j\leq n}\left(  x_{i}-x_{j}\right)  $,
we obtain a sum of several (in fact, $2^{n\left(  n-1\right)  /2}$ many)
monomials with $+$ and $-$ signs. I claim that among these monomials, the
monomial $x_{1}^{n-1}x_{2}^{n-2}\cdots x_{n}^{n-n}$ will appear exactly once,
and with a $+$ sign. Indeed, in order to obtain $x_{1}^{n-1}x_{2}^{n-2}\cdots
x_{n}^{n-n}$ when expanding the product
\[%
\begin{array}
[c]{ccccc}%
\prod_{1\leq i<j\leq n}\left(  x_{i}-x_{j}\right)  = & \left(  x_{1}%
-x_{2}\right)  & \left(  x_{1}-x_{3}\right)  & \cdots & \left(  x_{1}%
-x_{n}\right) \\
&  & \left(  x_{2}-x_{3}\right)  & \cdots & \left(  x_{2}-x_{n}\right) \\
&  &  & \ddots & \vdots\\
&  &  &  & \left(  x_{n-1}-x_{n}\right)  ,
\end{array}
\]
it is necessary to pick the $x_{1}$ minuends from all $n-1$ factors in the
first row (since none of the other factors contain any $x_{1}$), then to pick
the $x_{2}$ minuends from all $n-2$ factors in the second row (since none of
the remaining factors contain any $x_{2}$), and so on -- i.e., to take the
minuend (rather than the subtrahend) from each factor. Thus, only one of the
monomials obtained by the expansion will be $x_{1}^{n-1}x_{2}^{n-2}\cdots
x_{n}^{n-n}$, and it will appear with a $+$ sign. Hence, the coefficient of
$x_{1}^{n-1}x_{2}^{n-2}\cdots x_{n}^{n-n}$ in the product $\prod_{1\leq
i<j\leq n}\left(  x_{i}-x_{j}\right)  $ is $1$. In other words, the
coefficient of $x_{1}^{n-1}x_{2}^{n-2}\cdots x_{n}^{n-n}$ in $g$ is $1$ (since
$g=\prod_{1\leq i<j\leq n}\left(  x_{i}-x_{j}\right)  $). In other words,%
\[
\left[  x_{1}^{n-1}x_{2}^{n-2}\cdots x_{n}^{n-n}\right]  g=1.
\]


Now, recall that $f=gq$. Hence,%
\begin{align*}
\left[  x_{1}^{n-1}x_{2}^{n-2}\cdots x_{n}^{n-n}\right]  f  &  =\left[
x_{1}^{n-1}x_{2}^{n-2}\cdots x_{n}^{n-n}\right]  \left(  gq\right) \\
&  =q\cdot\underbrace{\left[  x_{1}^{n-1}x_{2}^{n-2}\cdots x_{n}^{n-n}\right]
g}_{=1}\ \ \ \ \ \ \ \ \ \ \left(  \text{since }q\in\mathbb{Z}\right) \\
&  =q,
\end{align*}
so that $q=\left[  x_{1}^{n-1}x_{2}^{n-2}\cdots x_{n}^{n-n}\right]  f=1$ (by
(\ref{pf.lem.det.vander.a.pol.f-LC})). Hence, $f=g\underbrace{q}_{=1}=g$. As
we said, this completes the proof of Lemma \ref{lem.det.vander.a.pol}.
\end{proof}

The technique used in the above proof of Lemma \ref{lem.det.vander.a.pol} may
appear somewhat underhanded: Instead of computing our polynomial $f$ upfront,
we have kept lopping off linear factors from it until a constant polynomial
remained (for degree reasons). This technique is called \emph{identification
of factors} or \emph{factor hunting}, and is used in various different places,
but particularly often in the computation of determinants (multiple examples
are given in \cite[\S 2.4 and further below]{Krattenthaler}). While I consider
it to be aesthetically inferior to sufficiently direct approaches, it has
shown to be useful in situations in which no direct approaches are known to work.

\begin{proof}
[Proof of Theorem \ref{thm.det.vander} (sketched).]\textbf{(a)} We have
already given a proof of Theorem \ref{thm.det.vander} \textbf{(a)} (and with
Lemma \ref{lem.det.vander.a.pol} established, this proof is now complete).
\medskip

\textbf{(b)} The matrix $\left(  a_{j}^{n-i}\right)  _{1\leq i\leq n,\ 1\leq
j\leq n}$ is the transpose of the matrix $\left(  a_{i}^{n-j}\right)  _{1\leq
i\leq n,\ 1\leq j\leq n}$. Thus, Theorem \ref{thm.det.vander} \textbf{(b)}
follows from Theorem \ref{thm.det.vander} \textbf{(a)} using Theorem
\ref{thm.det.transp}. \medskip

\textbf{(c)} Let $A$ be the $n\times n$-matrix $\left(  a_{i}^{n-j}\right)
_{1\leq i\leq n,\ 1\leq j\leq n}$. Then, Theorem \ref{thm.det.vander}
\textbf{(a)} says that $\det A=\prod_{1\leq i<j\leq n}\left(  a_{i}%
-a_{j}\right)  $.

Let $\tau\in S_{n}$ be the permutation of $\left[  n\right]  $ that sends each
$j\in\left[  n\right]  $ to $n+1-j$. Then, each $i,j\in\left[  n\right]  $
satisfy $j-1=n-\tau\left(  j\right)  $ and therefore
\[
a_{i}^{j-1}=a_{i}^{n-\tau\left(  j\right)  }=A_{i,\tau\left(  j\right)
}\ \ \ \ \ \ \ \ \ \ \left(  \text{by the definition of }A\right)  .
\]
Hence, $\left(  a_{i}^{j-1}\right)  _{1\leq i\leq n,\ 1\leq j\leq n}=\left(
A_{i,\tau\left(  j\right)  }\right)  _{1\leq i\leq n,\ 1\leq j\leq n}$ and
thus%
\[
\det\left(  \left(  a_{i}^{j-1}\right)  _{1\leq i\leq n,\ 1\leq j\leq
n}\right)  =\det\left(  \left(  A_{i,\tau\left(  j\right)  }\right)  _{1\leq
i\leq n,\ 1\leq j\leq n}\right)  =\left(  -1\right)  ^{\tau}\cdot\det A
\]
(by (\ref{eq.cor.det.sig-row-col.col})).

However, the permutation $\tau$ has OLN $\left(  n,n-1,n-2,\ldots,1\right)  $.
Thus, each pair $\left(  i,j\right)  \in\left[  n\right]  ^{2}$ satisfying
$i<j$ is an inversion of $\tau$. Therefore, the Coxeter length $\ell\left(
\tau\right)  $ of $\tau$ is the \# of all pairs $\left(  i,j\right)
\in\left[  n\right]  ^{2}$ satisfying $i<j$. Thus, the sign of $\tau$ is%
\begin{align*}
\left(  -1\right)  ^{\tau}  &  =\left(  -1\right)  ^{\ell\left(  \tau\right)
}\ \ \ \ \ \ \ \ \ \ \left(  \text{by the definition of the sign of a
permutation}\right) \\
&  =\left(  -1\right)  ^{\left(  \text{\# of all pairs }\left(  i,j\right)
\in\left[  n\right]  ^{2}\text{ satisfying }i<j\right)  }\\
&  \ \ \ \ \ \ \ \ \ \ \ \ \ \ \ \ \ \ \ \ \left(  \text{since }\ell\left(
\tau\right)  \text{ is the \# of all pairs }\left(  i,j\right)  \in\left[
n\right]  ^{2}\text{ satisfying }i<j\right) \\
&  =\prod_{1\leq i<j\leq n}\left(  -1\right)  .
\end{align*}


Now,
\begin{align*}
&  \det\left(  \left(  a_{i}^{j-1}\right)  _{1\leq i\leq n,\ 1\leq j\leq
n}\right) \\
&  =\underbrace{\left(  -1\right)  ^{\tau}}_{=\prod_{1\leq i<j\leq n}\left(
-1\right)  }\cdot\underbrace{\det A}_{=\prod_{1\leq i<j\leq n}\left(
a_{i}-a_{j}\right)  }=\left(  \prod_{1\leq i<j\leq n}\left(  -1\right)
\right)  \cdot\left(  \prod_{1\leq i<j\leq n}\left(  a_{i}-a_{j}\right)
\right) \\
&  =\prod_{1\leq i<j\leq n}\underbrace{\left(  \left(  -1\right)  \left(
a_{i}-a_{j}\right)  \right)  }_{=a_{j}-a_{i}}=\prod_{1\leq i<j\leq n}\left(
a_{j}-a_{i}\right)  =\prod_{1\leq j<i\leq n}\left(  a_{i}-a_{j}\right)
\end{align*}
(here, we have renamed the index $\left(  i,j\right)  $ as $\left(
j,i\right)  $ in the product). This proves Theorem \ref{thm.det.vander}
\textbf{(c)}. \medskip

\textbf{(d)} The matrix $\left(  a_{j}^{i-1}\right)  _{1\leq i\leq n,\ 1\leq
j\leq n}$ is the transpose of the matrix $\left(  a_{i}^{j-1}\right)  _{1\leq
i\leq n,\ 1\leq j\leq n}$. Thus, Theorem \ref{thm.det.vander} \textbf{(d)}
follows from Theorem \ref{thm.det.vander} \textbf{(c)} using Theorem
\ref{thm.det.transp}.
\end{proof}

The Vandermonde determinant is itself a useful tool in the computation of
various other determinants. Here is an example:

\begin{proposition}
\label{prop.det.(xi+yj)n-1}Let $n\in\mathbb{N}$. Let $x_{1},x_{2},\ldots
,x_{n}\in K$ and $y_{1},y_{2},\ldots,y_{n}\in K$. Then,%
\begin{align*}
&  \det\left(  \left(  \left(  x_{i}+y_{j}\right)  ^{n-1}\right)  _{1\leq
i\leq n,\ 1\leq j\leq n}\right) \\
&  =\det\left(
\begin{array}
[c]{cccc}%
\left(  x_{1}+y_{1}\right)  ^{n-1} & \left(  x_{1}+y_{2}\right)  ^{n-1} &
\cdots & \left(  x_{1}+y_{n}\right)  ^{n-1}\\
\left(  x_{2}+y_{1}\right)  ^{n-1} & \left(  x_{2}+y_{2}\right)  ^{n-1} &
\cdots & \left(  x_{2}+y_{n}\right)  ^{n-1}\\
\vdots & \vdots & \ddots & \vdots\\
\left(  x_{n}+y_{1}\right)  ^{n-1} & \left(  x_{n}+y_{2}\right)  ^{n-1} &
\cdots & \left(  x_{n}+y_{n}\right)  ^{n-1}%
\end{array}
\right) \\
&  =\left(  \prod_{k=0}^{n-1}\dbinom{n-1}{k}\right)  \left(  \prod_{1\leq
i<j\leq n}\left(  x_{i}-x_{j}\right)  \right)  \left(  \prod_{1\leq i<j\leq
n}\left(  y_{j}-y_{i}\right)  \right)  .
\end{align*}

\end{proposition}

\begin{proof}
[First proof of Proposition \ref{prop.det.(xi+yj)n-1} (sketched).]Here is a
rough outline of a proof that uses factor hunting (in the same way as in our
above proofs of Theorem \ref{thm.det.vander} \textbf{(a)} and Lemma
\ref{lem.det.vander.a.pol}). We WLOG assume that $x_{1},x_{2},\ldots,x_{n}$
and $y_{1},y_{2},\ldots,y_{n}$ are distinct indeterminates in a polynomial
ring over $\mathbb{Z}$. (This is again an assumption that we can make, because
the argument that we used to derive Theorem \ref{thm.det.vander} \textbf{(a)}
from Lemma \ref{lem.det.vander.a.pol} can be applied here as well.) Then, we
can easily see that%
\[
\det\left(  \left(  \left(  x_{i}+y_{j}\right)  ^{n-1}\right)  _{1\leq i\leq
n,\ 1\leq j\leq n}\right)
\]
is a homogeneous polynomial of degree $n\left(  n-1\right)  $. This polynomial
vanishes if we set any $x_{u}$ equal to any $x_{v}$ (for $u<v$), and also
vanishes if we set any $y_{u}$ equal to any $y_{v}$ (for $u<v$). Thus we have
identified $n\left(  n-1\right)  $ linear factors of this polynomial (namely,
the differences $x_{u}-x_{v}$ and $y_{v}-y_{u}$ for $u<v$), and we can again
conclude (since any polynomial ring over $\mathbb{Z}$ is a unique
factorization domain) that
\[
\det\left(  \left(  \left(  x_{i}+y_{j}\right)  ^{n-1}\right)  _{1\leq i\leq
n,\ 1\leq j\leq n}\right)  =\left(  \prod_{1\leq i<j\leq n}\left(  x_{i}%
-x_{j}\right)  \right)  \left(  \prod_{1\leq i<j\leq n}\left(  y_{j}%
-y_{i}\right)  \right)  \cdot q
\]
for a constant $q\in\mathbb{Z}$. It remains to prove that this constant $q$
equals $\prod_{k=0}^{n-1}\dbinom{n-1}{k}$. This can be done by studying the
coefficients of the monomial%
\[
x_{1}^{n-1}x_{2}^{n-2}\cdots x_{n}^{n-n}y_{1}^{0}y_{2}^{1}\cdots y_{n}^{n-1}%
\]
in $\det\left(  \left(  \left(  x_{i}+y_{j}\right)  ^{n-1}\right)  _{1\leq
i\leq n,\ 1\leq j\leq n}\right)  $ and in $\left(  \prod_{1\leq i<j\leq
n}\left(  x_{i}-x_{j}\right)  \right)  \left(  \prod_{1\leq i<j\leq n}\left(
y_{j}-y_{i}\right)  \right)  $. We leave the details to the reader.
\end{proof}

\begin{proof}
[Second proof of Proposition \ref{prop.det.(xi+yj)n-1}.](See \cite[Exercise
6.17 \textbf{(b)}]{detnotes} for details.) Let
\[
C:=\left(  \left(  x_{i}+y_{j}\right)  ^{n-1}\right)  _{1\leq i\leq n,\ 1\leq
j\leq n}\in K^{n\times n}.
\]
For any $i,j\in\left[  n\right]  $, we have%
\begin{align*}
C_{i,j}  &  =\left(  x_{i}+y_{j}\right)  ^{n-1}=\sum_{k=0}^{n-1}\dbinom
{n-1}{k}x_{i}^{k}y_{j}^{n-1-k}\ \ \ \ \ \ \ \ \ \ \left(  \text{by the
binomial formula}\right) \\
&  =\sum_{k=1}^{n}\dbinom{n-1}{k-1}x_{i}^{k-1}y_{j}^{n-k}%
\end{align*}
(here, we have substituted $k-1$ for $k$ in the sum). If we define two
$n\times n$-matrices $P$ and $Q$ by
\[
P:=\left(  \dbinom{n-1}{k-1}x_{i}^{k-1}\right)  _{1\leq i\leq n,\ 1\leq k\leq
n}\ \ \ \ \ \ \ \ \ \ \text{and}\ \ \ \ \ \ \ \ \ \ Q:=\left(  y_{j}%
^{n-k}\right)  _{1\leq k\leq n,\ 1\leq j\leq n},
\]
then we can rewrite this as%
\[
C_{i,j}=\sum_{k=1}^{n}P_{i,k}Q_{k,j}=\left(  PQ\right)  _{i,j}%
\]
(by the definition of the matrix product). Since this holds for all
$i,j\in\left[  n\right]  $, we thus obtain $C=PQ$. Hence,
\begin{equation}
\det C=\det\left(  PQ\right)  =\det P\cdot\det Q
\label{pf.prop.det.(xi+yj)n-1.2nd.detA=1}%
\end{equation}
(by Theorem \ref{thm.det.detAB}). It now remains to compute $\det P$ and $\det
Q$.

From $Q=\left(  y_{j}^{n-k}\right)  _{1\leq k\leq n,\ 1\leq j\leq n}=\left(
y_{j}^{n-i}\right)  _{1\leq i\leq n,\ 1\leq j\leq n}$, we obtain%
\begin{equation}
\det Q=\det\left(  \left(  y_{j}^{n-i}\right)  _{1\leq i\leq n,\ 1\leq j\leq
n}\right)  =\prod_{1\leq i<j\leq n}\left(  y_{i}-y_{j}\right)
\label{pf.prop.det.(xi+yj)n-1.2nd.detQ=}%
\end{equation}
(by Theorem \ref{thm.det.vander} \textbf{(b)}, applied to $a_{i}=y_{i}$).

From $P=\left(  \dbinom{n-1}{k-1}x_{i}^{k-1}\right)  _{1\leq i\leq n,\ 1\leq
k\leq n}=\left(  \dbinom{n-1}{j-1}x_{i}^{j-1}\right)  _{1\leq i\leq n,\ 1\leq
j\leq n}$, we obtain%
\begin{align}
\det P  &  =\det\left(  \left(  \dbinom{n-1}{j-1}x_{i}^{j-1}\right)  _{1\leq
i\leq n,\ 1\leq j\leq n}\right) \nonumber\\
&  =\underbrace{\dbinom{n-1}{1-1}\dbinom{n-1}{2-1}\cdots\dbinom{n-1}{n-1}%
}_{=\prod_{k=0}^{n-1}\dbinom{n-1}{k}}\cdot\underbrace{\det\left(  \left(
x_{i}^{j-1}\right)  _{1\leq i\leq n,\ 1\leq j\leq n}\right)  }%
_{\substack{=\prod_{1\leq j<i\leq n}\left(  x_{i}-x_{j}\right)  \\\text{(by
Theorem \ref{thm.det.vander} \textbf{(c)},}\\\text{applied to }a_{i}%
=x_{i}\text{)}}}\nonumber\\
&  \ \ \ \ \ \ \ \ \ \ \ \ \ \ \ \ \ \ \ \ \left(
\begin{array}
[c]{c}%
\text{by (\ref{eq.cor.det.scale-row-col.col}), applied to }A=\left(
x_{i}^{j-1}\right)  _{1\leq i\leq n,\ 1\leq j\leq n}\\
\text{and }d_{i}=\dbinom{n-1}{i-1}%
\end{array}
\right) \nonumber\\
&  =\left(  \prod_{k=0}^{n-1}\dbinom{n-1}{k}\right)  \cdot\prod_{1\leq j<i\leq
n}\left(  x_{i}-x_{j}\right) \nonumber\\
&  =\left(  \prod_{k=0}^{n-1}\dbinom{n-1}{k}\right)  \cdot\prod_{1\leq i<j\leq
n}\left(  x_{j}-x_{i}\right)  \label{pf.prop.det.(xi+yj)n-1.2nd.detP=}%
\end{align}
(here, we have renamed the index $\left(  j,i\right)  $ as $\left(
i,j\right)  $ in the second product).

Now, (\ref{pf.prop.det.(xi+yj)n-1.2nd.detA=1}) becomes%
\begin{align*}
\det C  &  =\det P\cdot\det Q\\
&  =\left(  \prod_{k=0}^{n-1}\dbinom{n-1}{k}\right)  \cdot\left(  \prod_{1\leq
i<j\leq n}\left(  x_{j}-x_{i}\right)  \right)  \cdot\prod_{1\leq i<j\leq
n}\left(  y_{i}-y_{j}\right) \\
&  \ \ \ \ \ \ \ \ \ \ \ \ \ \ \ \ \ \ \ \ \left(  \text{by
(\ref{pf.prop.det.(xi+yj)n-1.2nd.detP=}) and
(\ref{pf.prop.det.(xi+yj)n-1.2nd.detQ=})}\right) \\
&  =\left(  \prod_{k=0}^{n-1}\dbinom{n-1}{k}\right)  \cdot\prod_{1\leq i<j\leq
n}\underbrace{\left(  \left(  x_{j}-x_{i}\right)  \left(  y_{i}-y_{j}\right)
\right)  }_{=\left(  x_{i}-x_{j}\right)  \left(  y_{j}-y_{i}\right)  }\\
&  =\left(  \prod_{k=0}^{n-1}\dbinom{n-1}{k}\right)  \cdot\prod_{1\leq i<j\leq
n}\left(  \left(  x_{i}-x_{j}\right)  \left(  y_{j}-y_{i}\right)  \right) \\
&  =\left(  \prod_{k=0}^{n-1}\dbinom{n-1}{k}\right)  \left(  \prod_{1\leq
i<j\leq n}\left(  x_{i}-x_{j}\right)  \right)  \left(  \prod_{1\leq i<j\leq
n}\left(  y_{j}-y_{i}\right)  \right)  .
\end{align*}
This proves Proposition \ref{prop.det.(xi+yj)n-1} (since $C=\left(  \left(
x_{i}+y_{j}\right)  ^{n-1}\right)  _{1\leq i\leq n,\ 1\leq j\leq n}$).
\end{proof}

\subsubsection{Laplace expansion}

Let us next recall another fundamental property of determinants: \emph{Laplace
expansion}. We will use the following notation:

\begin{convention}
\label{conv.mat.tilde}Let $n\in\mathbb{N}$. Let $A$ be an $n\times n$-matrix.
Let $i,j\in\left[  n\right]  $. Then, we set%
\[
A_{\sim i,\sim j}:=\operatorname*{sub}\nolimits_{\left[  n\right]
\setminus\left\{  i\right\}  }^{\left[  n\right]  \setminus\left\{  j\right\}
}A\ \ \ \ \ \ \ \ \ \ \left(  \text{using the notation from Definition
\ref{def.det.sub}}\right)  .
\]
This is the $\left(  n-1\right)  \times\left(  n-1\right)  $-matrix obtained
from $A$ by removing its $i$-th row and its $j$-th column.
\end{convention}

\begin{example}
If $A=\left(
\begin{array}
[c]{ccc}%
a & b & c\\
a^{\prime} & b^{\prime} & c^{\prime}\\
a^{\prime\prime} & b^{\prime\prime} & c^{\prime\prime}%
\end{array}
\right)  $, then $A_{\sim1,\sim2}=\left(
\begin{array}
[c]{cc}%
a^{\prime} & c^{\prime}\\
a^{\prime\prime} & c^{\prime\prime}%
\end{array}
\right)  $.
\end{example}

Now, we can state the theorem underlying Laplace expansion:

\begin{theorem}
\label{thm.det.laplace}Let $n\in\mathbb{N}$. Let $A\in K^{n\times n}$ be an
$n\times n$-matrix. \medskip

\textbf{(a)} For every $p\in\left[  n\right]  $, we have%
\[
\det A=\sum_{q=1}^{n}\left(  -1\right)  ^{p+q}A_{p,q}\det\left(  A_{\sim
p,\sim q}\right)  .
\]


\textbf{(b)} For every $q\in\left[  n\right]  $, we have%
\[
\det A=\sum_{p=1}^{n}\left(  -1\right)  ^{p+q}A_{p,q}\det\left(  A_{\sim
p,\sim q}\right)  .
\]

\end{theorem}

Note that some authors denote the minors $\det\left(  A_{\sim p,\sim
q}\right)  $ in Theorem \ref{thm.det.laplace} by $A_{p,q}$. This is, of
course, totally incompatible with our notations.

\begin{proof}
[Proof of Theorem \ref{thm.det.laplace}.]See \cite[Theorem 6.82]{detnotes} or
\cite[Lemma 4.7.7]{Ford21} or \cite[5.8 and 5.8']{Laue-det} or
\cite[Proposition B.24 and Proposition B.25]{Strick13} or \cite[Theorem
9.48]{Loehr-BC}.
\end{proof}

Theorem \ref{thm.det.laplace} yields several ways to compute the determinant
of a matrix\footnote{In fact, some authors use Theorem \ref{thm.det.laplace}
as a \textbf{definition} of the determinant. (However, this is somewhat
tricky, as it requires proving that all the values obtained for $\det A$ by
applying Theorem \ref{thm.det.laplace} are actually equal.)}. When we compute
a determinant $\det A$ using Theorem \ref{thm.det.laplace} \textbf{(a)}, we
say that we \emph{expand this determinant along the }$p$\emph{-th row}. When
we compute a determinant $\det A$ using Theorem \ref{thm.det.laplace}
\textbf{(b)}, we say that we \emph{expand this determinant along the }%
$q$\emph{-th column}.

Theorem \ref{thm.det.laplace} has many applications, some of which you have
probably seen in your course on linear algebra. (A few might appear on the
homework.) The main theoretical application of Theorem \ref{thm.det.laplace}
is the concept of the \emph{adjugate} (or \emph{classical adjoint}) of a
matrix, which we shall introduce in a few moments.

First, let us see what happens if we replace the $A_{p,q}$ in Theorem
\ref{thm.det.laplace} by entries from another row or another column. In fact,
the respective sums become $0$ (instead of $\det A$), as the following
proposition shows:

\begin{proposition}
\label{prop.det.laplace.0}Let $n\in\mathbb{N}$. Let $A\in K^{n\times n}$ be an
$n\times n$-matrix. Let $r\in\left[  n\right]  $. \medskip

\textbf{(a)} For every $p\in\left[  n\right]  $ satisfying $p\neq r$, we have%
\[
0=\sum_{q=1}^{n}\left(  -1\right)  ^{p+q}A_{r,q}\det\left(  A_{\sim p,\sim
q}\right)  .
\]


\textbf{(b)} For every $q\in\left[  n\right]  $ satisfying $q\neq r$, we have%
\[
0=\sum_{p=1}^{n}\left(  -1\right)  ^{p+q}A_{p,r}\det\left(  A_{\sim p,\sim
q}\right)  .
\]

\end{proposition}

\begin{proof}
See \cite[Proposition 6.96]{detnotes}. This is also implicit in \cite[proof of
Proposition B.28]{Strick13} and \cite[proof of Theorem 9.50]{Loehr-BC}.
\medskip

\begin{fineprint}
Here is a sketch of the proof:

\textbf{(a)} Let $p\in\left[  n\right]  $ satisfy $p\neq r$. Let $C$ be the
matrix $A$ with its $p$-th row replaced by its $r$-th row. Then, the matrix
$C$ has two equal rows, so that $\det C=0$ (by Theorem \ref{thm.det.rowop}
\textbf{(c)}). On the other hand, expanding $\det C$ along the $p$-th row
(i.e., applying Theorem \ref{thm.det.laplace} \textbf{(a)} to $C$ instead of
$A$) yields%
\[
\det C=\sum_{q=1}^{n}\left(  -1\right)  ^{p+q}\underbrace{C_{p,q}}_{=A_{r,q}%
}\det\underbrace{\left(  C_{\sim p,\sim q}\right)  }_{=A_{\sim p,\sim q}}%
=\sum_{q=1}^{n}\left(  -1\right)  ^{p+q}A_{r,q}\det\left(  A_{\sim p,\sim
q}\right)  .
\]
Comparing these two equalities, we obtain the claim of Proposition
\ref{prop.det.laplace.0} \textbf{(a)}. A similar argument proves Proposition
\ref{prop.det.laplace.0} \textbf{(b)}.
\end{fineprint}
\end{proof}

We can now define the adjugate of a matrix:

\begin{definition}
\label{def.det.adj}Let $n\in\mathbb{N}$. Let $A\in K^{n\times n}$ be an
$n\times n$-matrix. We define a new $n\times n$-matrix $\operatorname*{adj}%
A\in K^{n\times n}$ by%
\[
\operatorname*{adj}A=\left(  \left(  -1\right)  ^{i+j}\det\left(  A_{\sim
j,\sim i}\right)  \right)  _{1\leq i\leq n,\ 1\leq j\leq n}.
\]
This matrix $\operatorname*{adj}A$ is called the \emph{adjugate} of the matrix
$A$. (Some older texts call it the \emph{adjoint}, but this name has since
been conquered by a different notion. As a compromise, some still call
$\operatorname*{adj}A$ the \emph{classical adjoint} of $A$.)
\end{definition}

\begin{example}
The adjugate of the $0\times0$-matrix is the $0\times0$-matrix.

The adjugate of a $1\times1$-matrix $\left(
\begin{array}
[c]{c}%
a
\end{array}
\right)  $ is $\operatorname*{adj}\left(
\begin{array}
[c]{c}%
a
\end{array}
\right)  =\left(
\begin{array}
[c]{c}%
1
\end{array}
\right)  $.

The adjugate of a $2\times2$-matrix $\left(
\begin{array}
[c]{cc}%
a & b\\
c & d
\end{array}
\right)  $ is
\[
\operatorname*{adj}\left(
\begin{array}
[c]{cc}%
a & b\\
c & d
\end{array}
\right)  =\left(
\begin{array}
[c]{cc}%
d & -b\\
-c & a
\end{array}
\right)  .
\]


The adjugate of a $3\times3$-matrix $\left(
\begin{array}
[c]{ccc}%
a & b & c\\
d & e & f\\
g & h & i
\end{array}
\right)  $ is
\[
\operatorname*{adj}\left(
\begin{array}
[c]{ccc}%
a & b & c\\
d & e & f\\
g & h & i
\end{array}
\right)  =\left(
\begin{array}
[c]{ccc}%
ei-fh & ch-bi & bf-ce\\
fg-di & ai-cg & cd-af\\
dh-ge & bg-ah & ae-bd
\end{array}
\right)  .
\]

\end{example}

The main property of the adjugate $\operatorname*{adj}A$ is its connection to
the inverse $A^{-1}$ of a matrix $A$. Indeed, if an $n\times n$-matrix $A$ is
invertible, then its inverse $A^{-1}$ is $\dfrac{1}{\det A}\cdot
\operatorname*{adj}A$. More generally, even if $A$ is not invertible, the
product of $\operatorname*{adj}A$ with $A$ (in either order) equals $\left(
\det A\right)  \cdot I_{n}$ (where $I_{n}$ is the $n\times n$ identity
matrix). Let us state this as a theorem:

\begin{theorem}
\label{thm.det.adj.inverse}Let $n\in\mathbb{N}$. Let $A\in K^{n\times n}$ be
an $n\times n$-matrix. Let $I_{n}$ denote the $n\times n$ identity matrix.
Then,%
\[
A\cdot\left(  \operatorname*{adj}A\right)  =\left(  \operatorname*{adj}%
A\right)  \cdot A=\left(  \det A\right)  \cdot I_{n}.
\]

\end{theorem}

\begin{proof}
See \cite[Theorem 6.100]{detnotes} or \cite[Lemma 4.7.9]{Ford21} or
\cite[Theorem 9.50]{Loehr-BC} or \cite[Proposition B.28]{Strick13}. \medskip

\begin{fineprint}
Here is a sketch of the argument: In order to show that $A\cdot\left(
\operatorname*{adj}A\right)  =\left(  \det A\right)  \cdot I_{n}$, it suffices
to check that the $\left(  i,j\right)  $-th entry of $A\cdot\left(
\operatorname*{adj}A\right)  $ equals $\det A$ whenever $i=j$ and equals $0$
otherwise. However, this follows from Theorem \ref{thm.det.laplace}
\textbf{(a)} (in the case $i=j$) and Proposition \ref{prop.det.laplace.0}
\textbf{(a)} (in the case $i\neq j$). Thus, $A\cdot\left(  \operatorname*{adj}%
A\right)  =\left(  \det A\right)  \cdot I_{n}$ is proved. Similarly, $\left(
\operatorname*{adj}A\right)  \cdot A=\left(  \det A\right)  \cdot I_{n}$ can
be shown.
\end{fineprint}
\end{proof}

More about the adjugate matrix can be found in \cite[\S 6.15]{detnotes} and
\cite[\S 5.4--\S 5.6]{trach}; see also \cite{Robins05} for some applications.

There is also a generalization of Theorem \ref{thm.det.laplace}, called
\emph{Laplace expansion along multiple rows (or columns)}:

\begin{theorem}
\label{thm.det.laplace-multi}Let $n\in\mathbb{N}$. Let $A\in K^{n\times n}$ be
an $n\times n$-matrix. We shall use the notations $\widetilde{I}$ and
$\operatorname*{sum}S$ as defined in Theorem \ref{thm.det.det(A+B)}. \medskip

\textbf{(a)} For every subset $P$ of $\left[  n\right]  $, we have%
\[
\det A=\sum_{\substack{Q\subseteq\left[  n\right]  ;\\\left\vert Q\right\vert
=\left\vert P\right\vert }}\left(  -1\right)  ^{\operatorname*{sum}%
P+\operatorname*{sum}Q}\det\left(  \operatorname*{sub}\nolimits_{P}%
^{Q}A\right)  \det\left(  \operatorname*{sub}\nolimits_{\widetilde{P}%
}^{\widetilde{Q}}A\right)  .
\]


\textbf{(b)} For every subset $Q$ of $\left[  n\right]  $, we have%
\[
\det A=\sum_{\substack{P\subseteq\left[  n\right]  ;\\\left\vert P\right\vert
=\left\vert Q\right\vert }}\left(  -1\right)  ^{\operatorname*{sum}%
P+\operatorname*{sum}Q}\det\left(  \operatorname*{sub}\nolimits_{P}%
^{Q}A\right)  \det\left(  \operatorname*{sub}\nolimits_{\widetilde{P}%
}^{\widetilde{Q}}A\right)  .
\]

\end{theorem}

\begin{example}
Let $n=4$ and $A=\left(
\begin{array}
[c]{cccc}%
a & b & c & d\\
a^{\prime} & b^{\prime} & c^{\prime} & d^{\prime}\\
a^{\prime\prime} & b^{\prime\prime} & c^{\prime\prime} & d^{\prime\prime}\\
a^{\prime\prime\prime} & b^{\prime\prime\prime} & c^{\prime\prime\prime} &
d^{\prime\prime\prime}%
\end{array}
\right)  $ and $P=\left\{  3,4\right\}  $. Then, Theorem
\ref{thm.det.laplace-multi} \textbf{(a)} says that%
\begin{align*}
&  \det A\\
&  =\sum_{\substack{Q\subseteq\left[  n\right]  ;\\\left\vert Q\right\vert
=\left\vert P\right\vert }}\left(  -1\right)  ^{\operatorname*{sum}%
P+\operatorname*{sum}Q}\det\left(  \operatorname*{sub}\nolimits_{P}%
^{Q}A\right)  \det\left(  \operatorname*{sub}\nolimits_{\widetilde{P}%
}^{\widetilde{Q}}A\right) \\
&  =\left(  -1\right)  ^{\operatorname*{sum}\left\{  3,4\right\}
+\operatorname*{sum}\left\{  1,2\right\}  }\det\left(  \operatorname*{sub}%
\nolimits_{\left\{  3,4\right\}  }^{\left\{  1,2\right\}  }A\right)
\det\left(  \operatorname*{sub}\nolimits_{\widetilde{\left\{  3,4\right\}  }%
}^{\widetilde{\left\{  1,2\right\}  }}A\right) \\
&  \ \ \ \ \ \ \ \ \ \ +\left(  -1\right)  ^{\operatorname*{sum}\left\{
3,4\right\}  +\operatorname*{sum}\left\{  1,3\right\}  }\det\left(
\operatorname*{sub}\nolimits_{\left\{  3,4\right\}  }^{\left\{  1,3\right\}
}A\right)  \det\left(  \operatorname*{sub}\nolimits_{\widetilde{\left\{
3,4\right\}  }}^{\widetilde{\left\{  1,3\right\}  }}A\right) \\
&  \ \ \ \ \ \ \ \ \ \ +\left(  -1\right)  ^{\operatorname*{sum}\left\{
3,4\right\}  +\operatorname*{sum}\left\{  1,4\right\}  }\det\left(
\operatorname*{sub}\nolimits_{\left\{  3,4\right\}  }^{\left\{  1,4\right\}
}A\right)  \det\left(  \operatorname*{sub}\nolimits_{\widetilde{\left\{
3,4\right\}  }}^{\widetilde{\left\{  1,4\right\}  }}A\right) \\
&  \ \ \ \ \ \ \ \ \ \ +\left(  -1\right)  ^{\operatorname*{sum}\left\{
3,4\right\}  +\operatorname*{sum}\left\{  2,3\right\}  }\det\left(
\operatorname*{sub}\nolimits_{\left\{  3,4\right\}  }^{\left\{  2,3\right\}
}A\right)  \det\left(  \operatorname*{sub}\nolimits_{\widetilde{\left\{
3,4\right\}  }}^{\widetilde{\left\{  2,3\right\}  }}A\right) \\
&  \ \ \ \ \ \ \ \ \ \ +\left(  -1\right)  ^{\operatorname*{sum}\left\{
3,4\right\}  +\operatorname*{sum}\left\{  2,4\right\}  }\det\left(
\operatorname*{sub}\nolimits_{\left\{  3,4\right\}  }^{\left\{  2,4\right\}
}A\right)  \det\left(  \operatorname*{sub}\nolimits_{\widetilde{\left\{
3,4\right\}  }}^{\widetilde{\left\{  2,4\right\}  }}A\right) \\
&  \ \ \ \ \ \ \ \ \ \ +\left(  -1\right)  ^{\operatorname*{sum}\left\{
3,4\right\}  +\operatorname*{sum}\left\{  3,4\right\}  }\det\left(
\operatorname*{sub}\nolimits_{\left\{  3,4\right\}  }^{\left\{  3,4\right\}
}A\right)  \det\left(  \operatorname*{sub}\nolimits_{\widetilde{\left\{
3,4\right\}  }}^{\widetilde{\left\{  3,4\right\}  }}A\right) \\
&  =\det\left(
\begin{array}
[c]{cc}%
a^{\prime\prime} & b^{\prime\prime}\\
a^{\prime\prime\prime} & b^{\prime\prime\prime}%
\end{array}
\right)  \det\left(
\begin{array}
[c]{cc}%
c & d\\
c^{\prime} & d^{\prime}%
\end{array}
\right)  -\det\left(
\begin{array}
[c]{cc}%
a^{\prime\prime} & c^{\prime\prime}\\
a^{\prime\prime\prime} & c^{\prime\prime\prime}%
\end{array}
\right)  \det\left(
\begin{array}
[c]{cc}%
b & d\\
b^{\prime} & d^{\prime}%
\end{array}
\right) \\
&  \ \ \ \ \ \ \ \ \ \ +\det\left(
\begin{array}
[c]{cc}%
a^{\prime\prime} & d^{\prime\prime}\\
a^{\prime\prime\prime} & d^{\prime\prime\prime}%
\end{array}
\right)  \det\left(
\begin{array}
[c]{cc}%
b & c\\
b^{\prime} & c^{\prime}%
\end{array}
\right)  +\det\left(
\begin{array}
[c]{cc}%
b^{\prime\prime} & c^{\prime\prime}\\
b^{\prime\prime\prime} & c^{\prime\prime\prime}%
\end{array}
\right)  \det\left(
\begin{array}
[c]{cc}%
a & d\\
a^{\prime} & d^{\prime}%
\end{array}
\right) \\
&  \ \ \ \ \ \ \ \ \ \ -\det\left(
\begin{array}
[c]{cc}%
b^{\prime\prime} & d^{\prime\prime}\\
b^{\prime\prime\prime} & d^{\prime\prime\prime}%
\end{array}
\right)  \det\left(
\begin{array}
[c]{cc}%
a & c\\
a^{\prime} & c^{\prime}%
\end{array}
\right)  +\det\left(
\begin{array}
[c]{cc}%
c^{\prime\prime} & d^{\prime\prime}\\
c^{\prime\prime\prime} & d^{\prime\prime\prime}%
\end{array}
\right)  \det\left(
\begin{array}
[c]{cc}%
a & b\\
a^{\prime} & b^{\prime}%
\end{array}
\right)  .
\end{align*}

\end{example}

\begin{proof}
[Proof of Theorem \ref{thm.det.laplace-multi}.]See \cite[Theorem
6.156]{detnotes}.
\end{proof}

\subsubsection{Desnanot--Jacobi and Dodgson condensation}

We come to more exotic results. The following theorem is one of several
versions of the \emph{Desnanot--Jacobi formula}:

\begin{theorem}
[Desnanot--Jacobi formula, take 1]\label{thm.det.des-jac-1}Let $n\in
\mathbb{N}$ be such that $n\geq2$. Let $A\in K^{n\times n}$ be an $n\times n$-matrix.

Let $A^{\prime}$ be the $\left(  n-2\right)  \times\left(  n-2\right)
$-matrix%
\[
\operatorname*{sub}\nolimits_{\left\{  2,3,\ldots,n-1\right\}  }^{\left\{
2,3,\ldots,n-1\right\}  }A=\left(  A_{i+1,j+1}\right)  _{1\leq i\leq
n-2,\ 1\leq j\leq n-2}.
\]
(This is precisely what remains of the matrix $A$ when we remove the first
row, the last row, the first column and the last column.) Then,%
\begin{align*}
\det A\cdot\det\left(  A^{\prime}\right)   &  =\det\left(  A_{\sim1,\sim
1}\right)  \cdot\det\left(  A_{\sim n,\sim n}\right)  -\det\left(
A_{\sim1,\sim n}\right)  \cdot\det\left(  A_{\sim n,\sim1}\right) \\
&  =\det\left(
\begin{array}
[c]{cc}%
\det\left(  A_{\sim1,\sim1}\right)  & \det\left(  A_{\sim1,\sim n}\right) \\
\det\left(  A_{\sim n,\sim1}\right)  & \det\left(  A_{\sim n,\sim n}\right)
\end{array}
\right)  .
\end{align*}

\end{theorem}

\begin{example}
For $n=3$, this is saying that%
\begin{align*}
&  \det\left(
\begin{array}
[c]{ccc}%
a & b & c\\
a^{\prime} & b^{\prime} & c^{\prime}\\
a^{\prime\prime} & b^{\prime\prime} & c^{\prime\prime}%
\end{array}
\right)  \cdot\det\left(
\begin{array}
[c]{c}%
b^{\prime}%
\end{array}
\right) \\
&  =\det\left(
\begin{array}
[c]{cc}%
b^{\prime} & c^{\prime}\\
b^{\prime\prime} & c^{\prime\prime}%
\end{array}
\right)  \cdot\det\left(
\begin{array}
[c]{cc}%
a & b\\
a^{\prime} & b^{\prime}%
\end{array}
\right)  -\det\left(
\begin{array}
[c]{cc}%
a^{\prime} & b^{\prime}\\
a^{\prime\prime} & b^{\prime\prime}%
\end{array}
\right)  \cdot\det\left(
\begin{array}
[c]{cc}%
b & c\\
b^{\prime} & c^{\prime}%
\end{array}
\right)  .
\end{align*}

\end{example}

\begin{proof}
[Proof of Theorem \ref{thm.det.des-jac-1}.]See \cite[Proposition
6.122]{detnotes} or \cite[\S 3.5, proof of Theorem 3.12]{Bresso99} or
\cite{zeilberger-twotime}.
\end{proof}

Theorem \ref{thm.det.des-jac-1} provides a recursive way of computing
determinants: Indeed, in the setting of Theorem \ref{thm.det.des-jac-1}, if
$\det\left(  A^{\prime}\right)  $ is invertible (which, when $K$ is a field,
simply means that $\det\left(  A^{\prime}\right)  \neq0$), then Theorem
\ref{thm.det.des-jac-1}) yields%
\begin{equation}
\det A=\dfrac{\det\left(  A_{\sim1,\sim1}\right)  \cdot\det\left(  A_{\sim
n,\sim n}\right)  -\det\left(  A_{\sim1,\sim n}\right)  \cdot\det\left(
A_{\sim n,\sim1}\right)  }{\det\left(  A^{\prime}\right)  }.
\label{eq.thm.det.des-jac-1.frac}%
\end{equation}
The five matrices appearing on the right hand side of this are smaller than
$A$, so their determinants are often easier to compute than $\det A$. In
particular, if you are proving something by strong induction on $n$, you will
occasionally be able to use the induction hypothesis to compute these
determinants. This method of recursively simplifying determinants is often
known as \emph{Dodgson condensation}, as it was popularized (perhaps even
discovered) by Charles Lutwidge Dodgson (aka Lewis Carroll) in \cite[Appendix
II]{Dodgso67}. We outline a sample application:

\begin{theorem}
[Cauchy determinant]\label{thm.det.cauchy}Let $n\in\mathbb{N}$. Let
$x_{1},x_{2},\ldots,x_{n}$ be $n$ elements of $K$. Let $y_{1},y_{2}%
,\ldots,y_{n}$ be $n$ elements of $K$. Assume that $x_{i}+y_{j}$ is invertible
in $K$ for each $\left(  i,j\right)  \in\left[  n\right]  ^{2}$. Then,%
\[
\det\left(  \left(  \dfrac{1}{x_{i}+y_{j}}\right)  _{1\leq i\leq n,\ 1\leq
j\leq n}\right)  =\dfrac{\prod\limits_{1\leq i<j\leq n}\left(  \left(
x_{i}-x_{j}\right)  \left(  y_{i}-y_{j}\right)  \right)  }{\prod
\limits_{\left(  i,j\right)  \in\left[  n\right]  ^{2}}\left(  x_{i}%
+y_{j}\right)  }.
\]

\end{theorem}

Once again, there are many ways to prove this (see, e.g., \cite[Exercise
6.18]{detnotes}, \cite[\S 1.3]{Prasolov}, \cite[Theorem 2]{Gri-19.9}, or
\url{https://proofwiki.org/wiki/Value_of_Cauchy_Determinant} ). But using the
Desnanot--Jacobi identity, there is a rather straightforward proof of Theorem
\ref{thm.det.cauchy}. Indeed, if $A$ is a \emph{Cauchy matrix} (i.e., a matrix
of the form $\left(  \dfrac{1}{x_{i}+y_{j}}\right)  _{1\leq i\leq n,\ 1\leq
j\leq n}$), then so is each submatrix of $A$. Thus, if we proceed by strong
induction on $n$, we can use the induction hypothesis to compute all five
determinants on the right hand side of (\ref{eq.thm.det.des-jac-1.frac}). The
only difficulty is making sure that $\det\left(  A^{\prime}\right)  $ is
invertible. To achieve this, we again have to WLOG assume that our $x$'s and
$y$'s are indeterminates in a polynomial ring, and we have to rewrite the
claim of Theorem \ref{thm.det.cauchy} in the form%
\[
\det\left(  \left(  \prod_{k\neq j}\left(  x_{i}+y_{k}\right)  \right)
_{1\leq i\leq n,\ 1\leq j\leq n}\right)  =\prod\limits_{1\leq i<j\leq
n}\left(  \left(  x_{i}-x_{j}\right)  \left(  y_{i}-y_{j}\right)  \right)
\]
in order for both sides to actually be polynomials in $\mathbb{Z}\left[
x_{1},x_{2},\ldots,x_{n},y_{1},y_{2},\ldots,y_{n}\right]  $. The details are
left to the reader.

Theorem \ref{thm.det.des-jac-1} can be generalized significantly. Here is the
simplest generalization, in which the special role played by the first and
last rows and the first and last columns is instead given to any two rows and
any two columns:

\begin{theorem}
\label{thm.det.des-jac-2}Let $n\in\mathbb{N}$ be such that $n\geq2$. Let $p$,
$q$, $u$ and $v$ be four elements of $\left[  n\right]  $ such that $p<q$ and
$u<v$. Let $A$ be an $n\times n$-matrix. Then,%
\begin{align*}
&  \det A\cdot\det\left(  \operatorname*{sub}\nolimits_{\left[  n\right]
\setminus\left\{  p,q\right\}  }^{\left[  n\right]  \setminus\left\{
u,v\right\}  }A\right) \\
&  =\det\left(  A_{\sim p,\sim u}\right)  \cdot\det\left(  A_{\sim q,\sim
v}\right)  -\det\left(  A_{\sim p,\sim v}\right)  \cdot\det\left(  A_{\sim
q,\sim u}\right)  .
\end{align*}

\end{theorem}

\begin{proof}
See \cite[Theorem 6.126]{detnotes}.
\end{proof}

Even more generally, \emph{Jacobi's complementary minor theorem for adjugates}
(appearing, e.g., in \cite[Theorem 5.22]{trach}, or in equivalent forms in
\cite[Theorem 2.5.2]{Prasolov} and \cite[(13)]{BruSch83}) says the following:

\begin{theorem}
[Jacobi's complementary minor theorem for adjugates]%
\label{thm.det.jacobi-complement}Let $n\in\mathbb{N}$. For any subset $I$ of
$\left[  n\right]  $, we let $\widetilde{I}$ denote the complement $\left[
n\right]  \setminus I$ of $I$. Set $\operatorname*{sum}S=\sum_{s\in S}s$ for
any finite set $S$ of integers. (For example, $\operatorname*{sum}\left\{
2,4,5\right\}  =2+4+5=11$.)

Let $A\in K^{n\times n}$ be an $n\times n$-matrix. Let $P$ and $Q$ be two
subsets of $\left[  n\right]  $ such that $\left\vert P\right\vert =\left\vert
Q\right\vert \geq1$. Then,%
\[
\det\left(  \operatorname*{sub}\nolimits_{P}^{Q}\left(  \operatorname*{adj}%
A\right)  \right)  =\left(  -1\right)  ^{\operatorname*{sum}%
P+\operatorname*{sum}Q}\left(  \det A\right)  ^{\left\vert Q\right\vert
-1}\det\left(  \operatorname*{sub}\nolimits_{\widetilde{Q}}^{\widetilde{P}%
}A\right)  .
\]

\end{theorem}

Theorem \ref{thm.det.des-jac-2} is the particular case of Theorem
\ref{thm.det.jacobi-complement} for $P=\left\{  u,v\right\}  $ and $Q=\left\{
p,q\right\}  $.\ \ \ \ \footnote{Indeed, if we set $P=\left\{  u,v\right\}  $
and $Q=\left\{  p,q\right\}  $ in the situation of Theorem
\ref{thm.det.des-jac-2}, then%
\begin{align*}
\operatorname*{sub}\nolimits_{P}^{Q}\left(  \operatorname*{adj}A\right)   &
=\left(
\begin{array}
[c]{cc}%
\left(  -1\right)  ^{u+p}\det\left(  A_{\sim p,\sim u}\right)  & \left(
-1\right)  ^{u+q}\det\left(  A_{\sim q,\sim u}\right) \\
\left(  -1\right)  ^{v+p}\det\left(  A_{\sim p,\sim v}\right)  & \left(
-1\right)  ^{v+q}\det\left(  A_{\sim q,\sim v}\right)
\end{array}
\right)  \ \ \ \ \ \ \ \ \ \ \text{and}\\
\operatorname*{sub}\nolimits_{\widetilde{Q}}^{\widetilde{P}}A  &
=\operatorname*{sub}\nolimits_{\left[  n\right]  \setminus\left\{
p,q\right\}  }^{\left[  n\right]  \setminus\left\{  u,v\right\}  }A;
\end{align*}
thus, Theorem \ref{thm.det.jacobi-complement} is easily seen to boil down to
Theorem \ref{thm.det.des-jac-2} in this case (the powers of $-1$ all cancel).}
Theorem \ref{thm.det.des-jac-1} is, of course, the particular case of Theorem
\ref{thm.det.des-jac-2} for $p=1$ and $q=n$ and $u=1$ and $v=n$.

\begin{noncompile}
TODO: Here I should eventually add a proof of Theorem
\ref{thm.det.jacobi-complement} a la Prasolov. However, it needs the
determinant of a block-triangular matrix.
\end{noncompile}
